\chapter{Conclusions and Future Perspectives}

The computational reconstruction of peptide sequences directly from tandem mass spectrometry fragmentation data without reference genomes or proteome databases---\emph{de novo} peptide sequencing---is indispensable for modern molecular medicine, oncology, and biotechnology. For years, the field faced a persistent trade-off: either accept the slow, $\mathcal{O}(L)$ sequential latency and error propagation of autoregressive Transformers, or face the severe discretization errors of continuous normalizing flows. This thesis addresses this trade-off through the development, mathematical analysis, and empirical validation of \textbf{DFlowNovo}.

\section{Summary of Contributions}

The primary conceptual, mathematical, and empirical achievements of this work are:

\begin{enumerate}[(1)]
\item \textbf{Discrete Flow Formulation:} By formulating sequence generation as a continuous-time jump process on the discrete amino acid simplex via Continuous-Time Markov Chains (CTMC), DFlowNovo eliminates continuous-to-discrete rounding errors. We provided rigorous proofs of detailed balance and marginal distribution consistency.
\item \textbf{Exact Polynomial-Time Knapsack Guidance (\texttt{KnapsackDP}):} We developed a GPU-accelerated dynamic programming reachability filter that evaluates continuous precursor mass constraints in $\mathcal{O}(1)$ time per token on a quantized mass grid ($0.02\,\text{Da}$ resolution), building upon foundational subset-sum and spectrum graph dynamic programming principles \citep{Dancik1999, Chen2001}. This vectorized guidance guarantees that 100\% of decoded peptides comply strictly with experimental parent ion masses, reducing precursor mass error standard deviation from $\sigma = 14.8\,\text{ppm}$ to $\sigma = 1.45\,\text{ppm}$.
\item \textbf{State-of-the-Art Biological Generalization:} On the 104,163 full test spectra of the standardized Nine-Species benchmark, DFlowNovo achieves \textbf{68.28\% strict exact match} and \textbf{83.01\% amino acid F1}, outperforming both releases of InstaNovo (v1.2 at 65.48\% and v1.0 at 53.20\%) and Casanovo (48.10\%). At 80\% precision, DFlowNovo accepts \textbf{86,412 spectra} (+31,412 over InstaNovo v1.0), representing a 57.1\% increase in accepted biological identifications.
\item \textbf{High-Throughput Non-Autoregressive Latency:} By generating all sequence positions in parallel within a fixed budget of $K=20$ integration steps, DFlowNovo achieves decoding speeds of \textbf{227.1 to 255.8 spectra per second} on an NVIDIA GPU. This provides a \textbf{4.33$\times$ to 4.94$\times$ throughput advantage} over autoregressive beam search across standard peptide lengths, widening to over 9$\times$ on long peptides ($L \ge 25$).
\item \textbf{Physical and Chemical First-Principles Grounding:} We conducted an exhaustive biophysical error analysis, demonstrating that 30.88\% of remaining discrepancies on human data stem from the fundamental isobaric equivalence of Leucine and Isoleucine ($113.084064\,\text{Da}$) under collisional dissociation, which requires odd-electron radical fragmentation ($w$-ions) to break.
\end{enumerate}

\section{Methodological Limitations}

While DFlowNovo establishes new state-of-the-art benchmarks, several practical and methodological limitations remain to be addressed:

\begin{enumerate}[(i)]
\item \textbf{Complex Post-Translational Modifications (PTMs):} Although DFlowNovo's modular vocabulary supports canonical modifications (such as carbamidomethylation of cysteine, oxidation of methionine, and phosphorylation of serine, threonine, and tyrosine), combinatorial combinations of open-search variable modifications expand the reachability grid dimensions and require expanded dynamic programming state spaces.
\item \textbf{Low-Resolution Ion-Trap Spectra:} DFlowNovo was optimized for high-resolution Orbitrap instruments operating at mass tolerances $\le 20\,\text{ppm}$ \citep{Makarov2000, Zubarev2007}. Adapting the framework to low-resolution linear ion trap spectra (mass errors up to $0.5\,\text{Da}$) requires wider pooling kernels that increase candidate branching factors.
\end{enumerate}

\section{Future Research Directions}

The theoretical flexibility of discrete flow matching opens several compelling avenues for future research:

\subsection{Multi-Modal Dissociation Flow Matching}
The single greatest physical barrier identified in this thesis is the isobaric ambiguity between Leucine and Isoleucine. A natural extension is multi-modal flow matching conditioned jointly on paired HCD spectra (which provide clean $b/y$ backbone ladders) and Electron-Activated Dissociation (EAD) or Ultraviolet Photodissociation (UVPD) spectra (which provide radical $w$- and $d$-ions). By conditioning the flow on both dissociation modalities simultaneously, DFlowNovo can achieve 100\% isobaric disambiguation without algorithmic guessing.

\subsection{Orthogonal Analytical Dimensions (RT and CCS)}
Modern mass spectrometers increasingly incorporate liquid chromatography retention time (RT) and trapped ion mobility spectrometry (TIMS) collision cross section (CCS). Because retention time correlates with peptide hydrophobicity and CCS reflects the gas-phase three-dimensional conformation of the ion, incorporating RT and CCS predictors as auxiliary energy terms during flow matching will further constrain the candidate sequence space.

\subsection{End-to-End De Novo Antibody and Immunopeptidome Assembly}
Finally, integrating DFlowNovo into an end-to-end overlap-graph assembler will enable the direct \emph{de novo} sequencing of full-length monoclonal antibodies, circulating immunoglobulin repertoires, and tumor-specific neoantigens directly from patient biofluids without requiring patient genomic sequencing, accelerating the development of personalized immunotherapies.
